% Part: applied-modal-logics
% Chapter: epistemic-logic
% Section: properties-accessibility

\documentclass[../../../include/open-logic-section]{subfiles}

\begin{document}

\olfileid[tr]{aml}{el}{acc}

\olsection{Erişilebilirlik Bağıntıları ve Epistemik İlkeler}

Normal modal mantıklardaki çerçeve karşılıklılığı hakkında zaten bildiklerimiz
ışığında, karakteristik !!{formula}s{}in epistemik yorumlar altında nasıl
göründüğünü incelemek isteyebiliriz. Epistemik mantıkların tipik olarak S5
içinde yorumlandığını belirtmiştik. Şimdi çeşitli epistemik ilkelerin nasıl
temsil edildiğine bakalım ve bunların çeşitli çerçeve koşullarına nasıl
karşılık geldiğini düşünelim.

Normal modal mantıktan, farklı modal !!{formula}s{}in erişilebilirlik
bağıntılarının farklı özelliklerini nitelediğini hatırlayın. Bu tablo, belirli
epistemik ilkelere karşılık gelen birkaçını seçmektedir.

\begin{table}[t]
    \begin{tabular}{| p{.48\textwidth} || p{.48\textwidth} |}
      \hline
      {\emph{$R$ \dots ise}} & {\emph{\dots, $\mModel{M}$ içinde doğrudur:}} \\
      \hline \hline
        
      & $\Knows (p \lif q) \lif (\Knows p \lif \Knows q)$ 
      \hfill \newline{(Kapanış)} \\
      \hline
      \emph{yansımalı}: $\forall w Rww$  
      & $\Knows p \lif p$ \hfill \newline{(Doğruluk)} \\
      \hline
      \emph{geçişli}: \newline
      $\forall u \forall v \forall w ((Ruv \land Rvw) \lif Ruw)$ & 
      $\Knows p \lif \Knows \Knows p$ \hfill 
           \newline{(Olumlu İçebakış)} \\
      \hline 
      \emph{Öklidyen}: \newline
      $\forall w \forall u \forall v ((Rwu \land Rwv) \lif Ruv)$ &  
      $\lnot \Knows p \lif \Knows \neg \Knows p$ \hfill 
        \newline{(Olumsuz İçebakış)} \\
      \hline
    \end{tabular}
    \caption{Dört epistemik ilke.}
    \ollabel{tab:four}
  \end{table} 

$T$ aksiyomuna karşılık gelen doğruluk, bu ilkeler arasında çoğu kez en az
tartışmalı olanı sayılır; çünkü !!a{formula} biliniyorsa onun doğru olması
gerektiği iddiasını temsil eder. Kapanış ile Olumlu ve Olumsuz İçebakış ise
çok daha tartışmalıdır.

$K$ aksiyomuna karşılık gelen kapanış, bir öznenin bilgisinin içerme altında
kapalı olduğu düşüncesini temsil eder. Bu, bazı durumlarda akla yatkın
görünebilir. Örneğin Victoria'daysam Vancouver Adası'nda olduğumu biliyor
olabilirim. Tuhaf kuşkucu senaryoları bir yana bırakırsak Victoria'da
olduğumu da bilirim; bu durum, Vancouver Adası'nda olduğumu bildiğimi de
düşündürmelidir. Dolayısıyla burada bilgimin mantıksal kapanışı görece
sezgisel görünebilir. Öte yandan bilgimizin bütün sonuçlarını her zaman
düşünmeyiz; bu nedenle kapanış başka durumlarda daha az sezgisel sonuçlara
yol açabilir.

Bazen KK ilkesi diye bilinen Olumlu İçebakış, kimi zaman bir şeyi biliyorsam
onu bildiğimi de bilirim önermesi olarak ifade edilir. Bu, 4 aksiyomunun
epistemik karşılığıdır. Benzer biçimde olumsuz içebakış, bir şeyi
\emph{bilmiyorsam} onu bilmediğimi bildiğim önermesi olarak ifade edilir ve 5
aksiyomuna karşılık gelir. Gerçekte bildiğim bir şeyi bilip bilmediğimden emin
olmadığım görece sıradan örnekler, bu iki ilkeye de karşıörnek gibi görünür.


\end{document}
